% Version as of July 19, 2012
% Version as of October 8, 1993

\chapter{Introduction to \texorpdfstring{\MPI/}{MPI}}
\label{sec:intro}
\label{chap:intro}
\label{sec:intro-2}
\label{chap:intro-2}

\section{Overview and Goals}
\MPI/ (Message-Passing Interface) is a \emph{message-passing library
  interface specification}.  All parts of this definition are
significant.  \MPI/ addresses primarily the message-passing parallel
programming model, in which data is moved from the address space of
one process to that of another process through cooperative operations
on each process.  Extensions to the ``classical'' message-passing model
are provided in collective operations, remote-memory access
operations, dynamic process creation, and parallel I/O.  \MPI/ is a
\emph{specification}, not an implementation; there are multiple
implementations of \MPI/.  This specification is for a \emph{library
  interface}; \MPI/ is not a language, and all \MPI/ operations are
expressed as functions, subroutines, or methods, according to the
appropriate language bindings
that, for C and Fortran, are part of the \MPI/ standard.  The standard has been
defined through an open process by a community of parallel computing
vendors, computer scientists, and application developers.  The next
few sections provide an overview of the history of \MPI/'s development.

The main advantages of establishing a message-passing standard are portability
and ease of use. In a distributed memory communication environment in which
the higher level routines and/or abstractions are
built
upon lower level
message-passing routines, the benefits of standardization are particularly
apparent. Furthermore, the definition of a message-passing standard, such
as that proposed here, provides vendors with a clearly defined base set of
routines that they can implement efficiently, or in some cases for which they can provide hardware support, thereby enhancing scalability.

The goal of the Message-Passing Interface, simply stated, is to
develop a widely used
standard for writing message-passing programs.
As such the interface should
establish a practical, portable, efficient, and flexible standard for message
passing.

A complete list of goals follows.

\begin{itemize}

\item
Design an application programming interface (not necessarily for compilers
or a system implementation library).

\item
Allow efficient communication: Avoid memory-to-memory copying,
allow overlap of computation and communication, and offload to communication
co-processors, where available.

\item
Allow for implementations that can be used in a heterogeneous environment.

\item
Allow convenient
C and Fortran
bindings for the interface.

\item
Assume a reliable communication interface:
the user need not cope with communication failures.
Such failures are dealt with by the underlying communication subsystem.

\item
Define an interface that can be implemented on many
vendor's platforms, with no significant changes in the underlying
communication and system software.

\item
Semantics of the interface should be language independent.

\item
The interface should be designed to allow for thread safety.
\end{itemize}


\section{Background of \texorpdfstring{\MPIIDOTO/}{MPI-1.0}}
\label{sec:background:i}

\MPI/ sought to make use of the most attractive features
of a number of existing message-passing systems, rather than selecting one of
them and adopting it as the standard. Thus, \MPI/ was strongly influenced
by work at the IBM T. J. Watson Research Center
\cite{IBM-report1,IBM-report2}, Intel's NX/2 \cite{NX-2}, Express
\cite{express}, nCUBE's Vertex \cite{Vertex},
p4 \cite{p4manual,p4paper}, and
PARMACS \cite{parmacs1,parmacs2}.
Other important contributions have come
from Zipcode\mpitermindex{Zipcode} \cite{zipcode1,zipcode2}, Chimp\mpitermindex{Chimp} \cite{chimp1,chimp2}, PVM\mpitermindex{PVM}
\cite{PVM2,PVM1}, Chameleon\mpitermindex{Chameleon} \cite{chameleon-user-ref}, and PICL\mpitermindex{PICL} \cite{picl}.


The \MPI/ standardization effort involved about 60 people from 40 organizations
mainly from the United States and Europe. Most of the major vendors of
concurrent computers were involved in \MPI/, along with researchers from
universities, government laboratories, and industry. The standardization
process began with the Workshop on
Standards for Message-Passing in a Distributed Memory Environment,
sponsored by the Center for Research on Parallel Computing, held April 29--30,
1992, in Williamsburg, Virginia
\cite{walker92b}. At this workshop the basic features essential
to a standard message-passing interface were discussed, and a working group
established to continue the standardization process.

A preliminary draft proposal, known as \MPII/, was put forward by Dongarra,
Hempel, Hey, and Walker in November 1992, and a revised version was completed
in February 1993 \cite{mpi1}. \MPII/ embodied the main features that were
identified at the Williamsburg workshop as being necessary in a message
passing standard.  Since \MPII/ was primarily intended to promote discussion
and ``get the ball rolling,'' it focused mainly on point-to-point
communications.  \MPII/ brought to the forefront a number of important
standardization issues, but did not include any collective communication
routines and was not thread-safe\mpitermindex{threads!thread-safe}.

In November 1992, a
meeting of the \MPI/ working group was held in Minneapolis, at which it was
decided to place the standardization process on a more formal footing, and
to generally adopt the procedures and organization of the High Performance
Fortran Forum. Subcommittees were formed for the major
component areas of the standard, and an email discussion service established
for each. In addition, the goal of producing a draft \MPI/ standard by the Fall
of 1993 was set.
To achieve this goal the \MPI/ working group met every
6 weeks for two days throughout the first 9 months of 1993, and
presented the draft \MPI/ standard at the
Supercomputing 93 conference in November 1993. These meetings and the email
discussion together constituted the \MPI/ Forum, membership of which has been open
to all members of the high performance computing community.


\section{Background of \texorpdfstring{\MPIIDOTI/}{MPI-1.1}, \texorpdfstring{\MPIIDOTII/}{MPI-1.2}, and \texorpdfstring{\MPIIIDOTO/}{MPI-2.0}}
\label{sec:background}

Beginning in March 1995, the \MPI/ Forum began meeting to
consider corrections and extensions to the original \MPI/ standard
document~\cite{MPISTDORG}.  The first product of these deliberations was
Version 1.1 of the \MPI/ specification, released in June of 1995
\cite{MPISTD}
(see \\
\url{http://www.mpi-forum.org} for official \MPI/ document releases).
At that time, effort focused
in
five areas.
\begin{enumerate}
\item Further corrections and clarifications for the \mpiidoti/ document.
\item Additions to \mpiidoti/ that do not significantly change its types of
  functionality (new datatype constructors, language interoperability, etc.).
\item Completely new types of functionality (dynamic processes, one-sided
  communication, parallel I/O, etc.) that are what everyone thinks of as
  ``\MPIII/ functionality.''
\item Bindings for Fortran 90 and C++.
\MPIII/ specifies
  C++ bindings
  for both \mpii/ and \mpiii/ functions, and extensions to the Fortran 77
  binding of \mpii/ and \mpiii/ to handle Fortran 90 issues.
\item Discussions of areas in which the \MPI/ process and framework seem likely
  to be useful, but where more discussion and experience are needed before
  standardization (e.g.,
  zero-copy
  semantics on shared-memory machines, real-time
  specifications).
\end{enumerate}

Corrections and clarifications (items of type 1 in the above list)
were
collected in Chapter~3 of the \MPIII/ document:
``Version 1.2 of \MPI/.''
That
chapter also contains the function for
identifying the version number.  Additions to \mpiidoti/ (items of
types 2, 3, and 4 in the above list) are in the
remaining chapters of the \MPIII/ document,
and constitute the specification for
\mpiii/.
Items of type 5 in the
above list have been moved to a separate document, the ``\MPI/ Journal
of Development'' (JOD), and are not part of the \mpiii/ standard.

This structure makes it easy for users and implementors to understand
what level of \MPI/ compliance a given implementation has:
\begin{itemize}
\item \mpii/ compliance will mean compliance with \mpiidotiii/.
  This is a useful level of compliance.  It means that the
  implementation conforms to the clarifications of \mpiidoti/ function
  behavior given in Chapter~3 of the \MPIII/ document.  Some
  implementations may require changes to be \mpii/ compliant.
\item \mpiii/ compliance will mean compliance with all of
  \mpiiidoti/.
\item The \MPI/ Journal of Development is not part of the \MPI/ standard.
\end{itemize}


It is to be emphasized that forward compatibility is preserved.  That is,
a valid \mpiidoti/ program is both a valid
\mpiidotiii/
program and a valid
\mpiiidoti/
program, and a valid
\mpiidotiii/
program is a valid
\mpiiidoti/
program.



\section{Background of \texorpdfstring{\MPIIDOTIII/}{MPI-1.3} and \texorpdfstring{\MPIIIDOTI/}{MPI-2.1}}
\label{sec:background:ii:i}
After the release of \MPIIIDOTO/, the \MPI/ Forum kept working on
errata and clarifications for both standard documents (\MPIIDOTI/ and \MPIIIDOTO/).
The short document ``Errata for \MPIIDOTI/'' was released October 12, 1998.
On
July 5, 2001, a first ballot of errata and clarifications for \MPIIIDOTO/
was released, and a second ballot was voted on May 22, 2002.
Both votes were done
electronically.
Both ballots were combined into one document:
``Errata
for \MPIII/,'' May
15, 2002.
This errata process was then interrupted, but the
Forum and its e-mail reflectors kept working
on new requests for clarification.

Restarting regular work of the \MPI/ Forum was initiated in
three meetings,
at EuroPVM/MPI'06 in Bonn, at EuroPVM/MPI'07 in Paris, and
at SC'07 in Reno.
In December 2007,
a steering committee started the
organization of new \MPI/ Forum meetings at regular 8-weeks intervals.
At the January 14--16, 2008
meeting in Chicago, the \MPI/ Forum decided to combine
the existing and future \MPI/ documents to one document for
each version of the \MPI/ standard.
For technical and historical reasons, this series was started with \MPIIDOTIII/.
Additional Ballots~3 and~4
solved
old questions from the
errata list started in 1995
up to
new questions from the last years.
After all documents (\MPIIDOTI/, \MPIII/, Errata for \MPIIDOTI/ (Oct.~12, 1998),
and \MPIIIDOTI/ Ballots 1--4) were combined into one draft document,
for each chapter, a chapter author and review team were defined.
They cleaned up the document to achieve a consistent \MPIIIDOTI/ document.
The
final \MPIIIDOTI/ standard document
was
finished in June 2008, and finally released with a second vote
in September 2008 in the meeting at Dublin,
just
before EuroPVM/MPI'08.

\section{Background of \texorpdfstring{\MPIIIDOTII/}{MPI-2.2}}
\MPIIIDOTII/ is a minor update to the \MPIIIDOTI/ standard.  This version
addresses additional errors and ambiguities that were not corrected in
the \MPIIIDOTI/ standard as well as a small number of extensions to
\MPIIIDOTI/ that met the following criteria:
\begin{itemize}
\item Any correct \MPIIIDOTI/ program is a correct \MPIIIDOTII/ program.
\item Any extension must have significant benefit for users.
\item Any extension must not require significant implementation
  effort.  To that end, all such changes are accompanied by an open
  source implementation.
\end{itemize}
The discussions of \MPIIIDOTII/ proceeded concurrently with the \MPIIII/
discussions; in some cases, extensions were proposed for \MPIIIDOTII/ but
were later moved to \MPIIII/.

\section{Background of \texorpdfstring{\MPIIIIDOTO/}{MPI-3.0}}
\MPIIIIDOTO/ is a major update to the \MPI/ standard.  The updates include
the extension of collective operations to
include nonblocking versions, extensions to the one-sided operations, and
a new Fortran 2008 binding.  In addition, the deprecated C++ bindings
have been removed, as well as many of the deprecated routines and \MPI/
objects (such as the \mpidtype{MPI\_UB} datatype). Any valid \MPIIIDOTII/
program not using any of these removed \MPI/ procedures or objects is a valid
\MPIIIIDOTO/ program.

\section{Background of \texorpdfstring{\MPIIIIDOTI/}{MPI-3.1}}
\MPIIIIDOTI/ is a minor update to the \MPI/ standard.  Most of the updates
are corrections and clarifications to the standard,
especially for the Fortran bindings.  New functions added
include routines to manipulate \code{MPI\_Aint} values in a portable manner, nonblocking
collective I/O routines, and routines to get the index value by name for
\mpiskipfunc{MPI\_T} performance and control variables.
A general index was also added. Any valid \MPIIIIDOTO/ program is a valid
\MPIIIIDOTI/ program.

\section{Background of \texorpdfstring{\MPIIVDOTO/}{MPI-4.0}}
\MPIIVDOTO/ is a major update to the \MPI/ standard.
The largest changes are the addition of large-count versions of many
routines to address the limitations of using an \code{int}
or \code{INTEGER} for the count parameter, persistent collectives, partitioned
communications, an alternative way to initialize \MPI/, application
info assertions, and improvements to the definitions of error handling.
In addition, there are a number of smaller improvements and corrections.
Any
valid \MPIIIIDOTI/ program is a valid \MPIIVDOTO/ program with the exception of
semantic changes listed in Chapter~\ref{chap:semantic-changes}.

\section{Background of \texorpdfstring{\MPIIVDOTI/}{MPI-4.1}}
\MPIIVDOTI/ is a minor update to the \MPI/ standard.
It contains mostly corrections
and clarifications to the \MPIIVDOTO/ document.
Several routines, the
attribute key \mpiconst{MPI\_HOST}, and the \code{mpif.h} Fortran
include file are deprecated.
A new routine provides a way to inquire about the hardware on which
the \MPI/ program is running.
Any valid \MPIIVDOTO/ program is a valid \MPIIVDOTI/ program with the
exception of semantic changes listed in
Chapter~\ref{chap:semantic-changes}.

\section{Background of \texorpdfstring{\MPIVDOTO/}{MPI-5.0}}
\MPIVDOTO/ is a major update to the \MPI/ standard.
The largest change is the addition of a standard Application
Binary Interface (ABI) to allow interoperability of different
implementations.
In addition, there are a number of smaller improvements and corrections.
Any valid \MPIIVDOTI/ program is a valid \MPIVDOTO/ program.

%% draft items will be removed after 5.1 is final
\ifdraftspec
\section{Background of 202X Draft Specification}
The 202X draft specification is expected to become the \MPIVDOTI/
specification once all features have been merged.
\else
\fi


\section{Who Should Use This Standard?}

This standard is intended for use by all those who want to write portable
message-passing programs in
Fortran and C (and access the C bindings
from C++).
This includes individual application programmers,
developers of software designed to run on parallel machines, and creators of
environments and tools.  In order to be
attractive to this wide audience, the standard must provide a simple, easy-to-use
interface for the basic user while not semantically precluding the
high-performance message-passing operations available on advanced machines.


\section{What Platforms Are Targets for Implementation?}

The attractiveness of the message-passing paradigm at least partially
stems from its wide portability.  Programs expressed this way may run
on distributed-memory
multiprocessors, networks of workstations, and combinations of all of these.
In addition, shared-memory implementations,
including those for multi-core processors and hybrid
architectures,
are possible.
The paradigm will not be made obsolete by architectures combining the shared-
and distributed-memory views, or by increases in network speeds.  It thus
should be both possible and useful to implement this standard on a great
variety of machines, including those ``machines'' consisting of collections of
other machines, parallel or not, connected by a communication network.

The interface is suitable for use by fully general MIMD (Multiple Instruction,
Multiple Data) programs, as well as
those written in the more restricted style of
SPMD (Single Program, Multiple Data).
\MPI/ provides many features intended to improve performance on
scalable parallel computers with
specialized interprocessor communication
hardware.  Thus, we expect that native, high-performance
implementations of \MPI/ will be provided on such machines.  At the
same time, implementations of \MPI/ on top of standard Unix
interprocessor communication protocols will provide portability to
workstation clusters and heterogenous networks of workstations.

\section{What Is Included in the Standard?}

The standard includes:

\begin{itemize}
    \item Point-to-point communication,
    \item Partitioned communication,
    \item Datatypes,
    \item Collective operations,
    \item Process groups,
    \item Communication contexts,
    \item Virtual Topologies for \MPI/ Processes,
    \item Environmental management and inquiry,
    \item The Info object,
    \item Process initialization, creation, and management,
    \item One-sided communication,
    \item External interfaces,
    \item Parallel file I/O,
    \item Tool support,
    \item Process fault tolerance,
    \item Language bindings for Fortran and C,
    \item Application Binary Interface, and
    \item Additional topics in side-documents.
\end{itemize}

\section{Side-documents}
\label{sec:side-documents}

Side-documents extend and/or modify features, semantics, language bindings, and other aspects covered in this document.
Side-documents shall not modify any aspects defined in the \MPI/ Standard without providing a mechanism that explicitly enables these deviations.
Execution of a program that does not explicitly enable deviations from the \MPI/ Standard will comply with the \MPI/ Standard, even when using an \MPI/ implementation that implements a side-document that modifies any aspects.

Each side-document is versioned with a scheme that is independent from the \MPI/ Standard version and from other side-documents.
All side-documents specify compatibility and interoperability with versions of the \MPI/ Standard and may define interoperability with features and semantics from other side-documents.
Side-documents are not required to provide full coverage of all \MPI/ concepts, but shall document which \MPI/ concepts are affected.
A compliant implementation is not required to comply with any side-documents.
However, if compliance with a particular version of a side-document is claimed, the implementation must comply with the entire side-document.
Side-documents will be found at the same location as the \MPI/ Standard~\cite{mpi-standard-location}.

\section{Organization of This Document}

The following is a list of the remaining chapters in this document, along with
a brief description of each.

\begin{itemize}
\item
Chapter~\ref{sec:terms}, \nameref{sec:terms},
explains notational terms and conventions used
throughout the \MPI/ document.
\item
Chapter~\ref{sec:pt2pt}, \nameref{sec:pt2pt},
defines the basic, pairwise communication subset of \MPI/.
\emph{Send}
and \emph{receive} are found here, along with many associated functions
designed to make basic communication powerful and efficient.
\item
Chapter~\ref{sec:part}, \nameref{sec:part},
defines a method of performing partitioned communication in \MPI/. 
Partitioned communication allows multiple contributions of data to be made, 
potentially, from multiple actors (e.g., threads or tasks) in an \MPI/ process 
to a single message.
\item
Chapter~\ref{chap:datatypes}, \nameref{chap:datatypes}, defines a method to describe any data layout,
e.g., an array of structures.
\item
Chapter~\ref{sec:coll}, \nameref{sec:coll}, defines
process-group collective communication operations.  Well known
examples of this are barrier and broadcast over a group of processes (not
necessarily all the processes).
\item
Chapter~\ref{sec:context}, \nameref{sec:context},
shows how groups of processes are formed and manipulated, how unique
communication contexts are obtained, and how the two are bound together
into a \emph{communicator}.
\item
Chapter~\ref{sec:topol}, \nameref{sec:topol}, explains a set
of utility functions meant to assist in the mapping of \MPI/ process
groups (a linearly ordered set) to richer topological structures
such as multi-dimensional grids.
\item
Chapter~\ref{sec:environment}, \nameref{sec:environment}, explains
how the programmer can manage and make inquiries of the current \MPI/ environment.
These functions are needed for the writing of correct, robust programs,
and are especially important for the construction of highly-portable
message-passing programs.
\end{itemize}
\begin{itemize}
\item
Chapter~\ref{sec:misc-2}, \nameref{sec:misc-2}, defines an opaque object that is used as input in several \MPI/ routines.
\item
Chapter~\ref{sec:dynamic-2}, \nameref{sec:dynamic-2},
defines
several approaches to \MPI/ initialization, process creation,
and process management while placing minimal restrictions on the execution environment.
\item
Chapter~\ref{sec:one-side-2}, \nameref{sec:one-side-2}, defines
communication routines that can be completed by a single process.  These include
shared-memory operations (put/get) and remote accumulate operations.
\item
Chapter~\ref{sec:ei-2}, \nameref{sec:ei-2},
defines routines designed to allow developers to layer on
top of \MPI/.
\item
Chapter~\ref{sec:io-2}, \nameref{sec:io-2}, defines
\MPI/
support for parallel I/O.
\item
Chapter~\ref{chap:tools}, \nameref{chap:tools},
covers interfaces that allow debuggers, performance
analyzers, and other tools to obtain data about the operation of \MPI/
processes.
\item
Chapter~\ref{chap:ft}, \nameref{chap:ft},
covers interfaces that allow developers to design applications tolerant to
\MPI/ process failures. The definitions and interfaces presented in this
chapter describe the behavior of \MPI/ operations after a process failure
error has been raised, and provide supplementary interfaces to restore the
communication capabilities of \MPI/.
\item
Chapter~\ref{chap:deprecated}, \nameref{chap:deprecated}, describes routines that
are kept for reference. However usage of these functions is discouraged, as
they may be deleted in future versions of the standard.
\item
Chapter~\ref{chap:removed}, \nameref{chap:removed}, describes
routines and constructs that have been removed from \MPI/.
\item
Chapter~\ref{chap:semantic-changes}, \nameref{chap:semantic-changes}, describes
semantic changes from previous versions of \MPI/.
\item
Chapter~\ref{sec:binding-2},
\nameref{sec:binding-2}, discusses Fortran issues
and describes language interoperability aspects between
C and Fortran.
\item
Chapter~\ref{sec:abi},
\nameref{sec:abi}, discusses Application Binary Interface issues
and describes interoperability of compiled code.
\end{itemize}

The Appendices are:

\begin{itemize}
\item
Annex~\ref{sec:lang},
\nameref{sec:lang},
gives specific syntax in
C and Fortran
for all \MPI/ functions,
constants, and types.

\item
Annex~\ref{chap:change},
\nameref{chap:change}, summarizes some changes since the previous
version of the standard.
\item
Several \textTitleFont{index} pages show the locations of
\hyperref[index:general]{general terms and definitions},    %%ALLOWLATEX%
\hyperref[index:example]{examples},    %%ALLOWLATEX%
\hyperref[index:constant]{constants and predefined handles}, %%ALLOWLATEX%
\hyperref[index:declaration]{declarations of C and Fortran types}, %%ALLOWLATEX%
\hyperref[index:callback]{callback routine prototypes}, and all %%ALLOWLATEX%
\hyperref[index:function]{\MPI/ functions}. %%ALLOWLATEX%

\end{itemize}

\mpi/
provides various interfaces to facilitate interoperability
  of distinct  \mpi/ implementations.  Among these are the canonical
  data representation for \mpi/ I/O and for \mpifunc{MPI\_PACK\_EXTERNAL} and
  \mpifunc{MPI\_UNPACK\_EXTERNAL}.
The definition of an actual binding of these interfaces that will
enable interoperability is outside the scope of this document.

A separate document consists of ideas that were discussed in the
\MPI/ Forum during the \MPIII/ development and
deemed to have value, but were not included in the \mpi/
standard.  They are part of the ``Journal of Development'' (JOD),
which was created to capture these ideas and discussions.
The JOD is available at \url{https://www.mpi-forum.org/docs}.

